\begin{figure}[ht]
\centering
\begin{tikzpicture}[x=1.0cm, y=1.0cm, font=\small]

% Axes (complex plane)
\draw[->, thick] (-1.0, 0) -- (9.5, 0) node[anchor=west, font=\scriptsize] {real};
\draw[->, thick] (4.0, -3.0) -- (4.0, 3.0) node[anchor=south, font=\scriptsize] {imag};

% Disc centre 2 (shifted to x=2 on plot from origin at x=4 => plot x=2? we place plot origin at 4)
% Map: plot_x = 4 + value. Centres at 2, 5, 8 -> plot 6, 9? keep small. Use centres 2,5 and radii 1, 1.5
% disc 1: centre 2 (plot 4+? ) simpler: put real origin at plot x=4, unit=1cm
% centre a11 = 2 -> plot (6,0), radius 1
\draw[thick, engblue, fill=engblue!12, fill opacity=0.5] (6, 0) circle (1);
\node[engblue, font=\scriptsize] at (6, 1.3) {$D_{1}$};
\fill[engblue] (6, 0) circle (1.4pt);
\node[engblue, font=\scriptsize, anchor=north] at (6, -0.15) {$2$};

% centre a22 = -1 -> plot (3,0), radius 0.5
\draw[thick, engred, fill=engred!12, fill opacity=0.5] (3, 0) circle (0.5);
\node[engred, font=\scriptsize] at (3, 0.85) {$D_{2}$};
\fill[engred] (3, 0) circle (1.4pt);
\node[engred, font=\scriptsize, anchor=north] at (3, -0.15) {$-1$};

% centre a33 = 5 -> plot (9,0) too far; use 4.5 -> plot (8.5,0) radius 0.8
\draw[thick, engamber, fill=engamber!12, fill opacity=0.5] (8.5, 0) circle (0.8);
\node[engamber, font=\scriptsize] at (8.5, 1.1) {$D_{3}$};
\fill[engamber] (8.5, 0) circle (1.4pt);
\node[engamber, font=\scriptsize, anchor=north] at (8.5, -0.15) {$4.5$};

% Real axis origin tick
\draw (4, 0.08) -- (4, -0.08) node[anchor=north east, font=\scriptsize] {$0$};

\node[anchor=north, font=\scriptsize, text width=11cm, align=center]
  at (4.0, -3.3) {
  Each Gershgorin disc is centred on a diagonal entry with radius the
  sum of the off-diagonal magnitudes in that row. Every eigenvalue
  lies in the union of the discs. No disc here contains the origin, so
  the matrix is nonsingular and its eigenvalues stay clear of zero, a
  read-off of invertibility that costs only a scan of the rows.
};

\end{tikzpicture}
\caption[Gershgorin discs localise the spectrum]{Gershgorin discs for
a $3 \times 3$ matrix, one disc per row, centred on the diagonal entry
$a_{ii}$ with radius $\sum_{j \neq i} \abs{a_{ij}}$. The theorem
guarantees every eigenvalue sits inside the union. When no disc
reaches the origin, as here, the matrix cannot have a zero eigenvalue
and is therefore invertible; when the discs are small and well
separated from zero, the matrix is well away from singularity. The
disc scan is the cheapest available bound on where a matrix's spectrum
can lie, and a strictly diagonally dominant matrix is exactly one
whose discs all miss the origin.}
\label{fig:vol02:ch09:gershgorin}
\end{figure}
